\documentclass[11pt]{article}
\usepackage[utf8]{inputenc}
\usepackage{amsmath,amssymb}
\usepackage{booktabs}
\usepackage{graphicx}
\usepackage{hyperref}
\usepackage[margin=1in]{geometry}
\usepackage{multirow}
\usepackage{xcolor}

\title{Architecture Determines Safety Monitoring Strategy:\\Hidden-State Detection in Reasoning Models\\Reveals Two Distinct Regimes}

\author{Anonymous}
\date{}

\begin{document}
\maketitle

\begin{abstract}
Reasoning models exhibit ``self-jailbreak'' behavior where safety reasoning and harmful execution interleave throughout extended chain-of-thought, culminating in a commitment point where safety reasoning permanently ceases.
We present the first cross-architecture evaluation of hidden-state safety monitoring for this phenomenon, testing across four reasoning models spanning two architecture families (Llama-3.1 and Qwen2.5) at two scales (7--8B and 32B).
Our key finding is that hidden-state safety monitoring effectiveness is \emph{architecture-dependent}: we identify two distinct regimes.
In the \textbf{distributed-signal regime} (Llama), safety signal is spread across minor principal components (PCA top-3 explains 16.6\% variance), requiring temporal context for reliable detection (+30.2pp improvement, $p<0.05$).
In the \textbf{concentrated-signal regime} (Qwen2.5), safety signal concentrates in dominant principal components (PCA top-3 explains 38--41\% variance), enabling static probes to achieve 70\%+ coverage with temporal context providing no statistically significant improvement (+4.6pp, CI crossing zero).
We propose a diagnostic framework using Fisher discriminant ratio and PCA variance concentration as predictors of which monitoring regime applies, demonstrate that cross-architecture transfer fails catastrophically (25--33\% coverage), and show that the concentrated-signal regime uniquely enables positive detection lead time.
These findings establish that practitioners must select monitoring strategies based on model architecture rather than applying universal approaches.
\end{abstract}

\section{Introduction}

Reasoning models such as DeepSeek-R1 and QwQ generate extended chains of thought (CoT) that interleave safety reasoning with task execution. Recent work has identified a ``self-jailbreak'' phenomenon: models may gradually cease safety-relevant reasoning during long CoT sequences, culminating in a \emph{commitment point} after which safety reasoning permanently disappears \cite{harmthoughts,selfjailbreak}.

Hidden-state monitoring offers a promising defense against this phenomenon, as it operates independently of potentially deceptive text-level outputs. Prior work on hidden-state jailbreak detection (TrajGuard, SafeSwitch, SIREN) has achieved high defense rates against adversarial inputs to instruction-tuned models, but these methods have not been evaluated on reasoning-model self-jailbreak where the threat originates from the model's own reasoning process.

We present the first systematic cross-architecture evaluation of hidden-state monitoring for safety commitment detection in reasoning models. Our study spans four models across two architecture families (Llama-3.1 and Qwen2.5) at two scales (7--8B and 32B), enabling us to disentangle the effects of architecture from scale.

Our central finding challenges the assumption that any single monitoring approach generalizes across architectures. We identify two distinct regimes:

\begin{enumerate}
    \item \textbf{Distributed-signal regime} (Llama-3.1): Safety signal is distributed across minor principal components. Static probes achieve only 46.5\% coverage. Temporal context is essential, yielding +30.2pp improvement (95\% CI: [+19, +42]).
    \item \textbf{Concentrated-signal regime} (Qwen2.5): Safety signal concentrates in top principal components. Static probes achieve 70\%+ coverage. Temporal context adds only +4.6pp (95\% CI: [-2, +12], not significant).
\end{enumerate}

Beyond this characterization, we contribute: (a) a diagnostic framework using Fisher discriminant ratio and PCA variance as predictors of monitoring regime; (b) evidence that architecture, not scale, determines the regime (via controlled 7B vs 32B comparisons within Qwen2.5); (c) demonstration that cross-architecture probe transfer fails catastrophically; and (d) the finding that the concentrated-signal regime uniquely enables positive pre-textual detection lead time.

\section{Related Work}

\paragraph{Hidden-state jailbreak detection.} TrajGuard monitors hidden-state trajectories via Mahalanobis distance for anomaly detection. SafeSwitch and SIREN train linear probes on hidden representations to classify safety-relevant states. These achieve 95\%+ defense rates on instruction-tuned models but target adversarial \emph{inputs}, not self-generated reasoning failures.

\paragraph{Reasoning model safety.} HarmThoughts provides the first characterization of safety reasoning patterns in extended CoT, identifying interleaving of safety and harmful content. Self-Jailbreaking demonstrates that reasoning models can circumvent their own safety training through extended deliberation. Neither work proposes deployable hidden-state detectors.

\paragraph{CoT monitoring.} SafeRemind and SafeThink propose text-level monitoring of reasoning chains. However, these are vulnerable to pathological CoT where the model reasons safely in text while encoding harmful intent in hidden states. Hidden-state monitoring provides a complementary, manipulation-resistant signal.

\section{Method}

\subsection{Task Formulation}

Given a reasoning trace $\mathbf{x} = (x_1, \ldots, x_T)$ with per-step safety labels $y_t \in \{0, 1\}$ (from the HarmThoughts 16-label taxonomy, binarized), we define the \emph{commitment point} $t^*$ as the last step where $y_t = 1$ (safety-relevant). The monitoring task is: detect $t^*$ from hidden states $\mathbf{h}_t \in \mathbb{R}^d$ extracted at each reasoning step.

\subsection{Detection Methods}

\paragraph{Static probe.} Logistic regression on a single hidden state: $\hat{y}_t = \sigma(\mathbf{w}^\top \mathbf{h}_t + b)$.

\paragraph{Temporal window probe.} Logistic regression on concatenated window of hidden states: $\hat{y}_t = \sigma(\mathbf{w}^\top [\mathbf{h}_{t-W+1}; \ldots; \mathbf{h}_t] + b)$ for window size $W$.

\paragraph{GRU sequential probe.} A GRU-256 processes the hidden state sequence and classifies at each step, providing an upper bound on temporal modeling capacity.

\paragraph{Mahalanobis distance baseline.} Following TrajGuard, we compute Mahalanobis distance from a fitted safe-state distribution and threshold for anomaly detection.

\subsection{Diagnostic Framework}

We propose two metrics to predict which monitoring regime applies to a given architecture:

\begin{itemize}
    \item \textbf{PCA variance concentration}: Fraction of variance explained by top-3 principal components of hidden states. High concentration ($>30\%$) predicts the concentrated-signal regime.
    \item \textbf{Fisher discriminant ratio (PCA-projected)}: Fisher ratio computed in the top-3 PC subspace. High PCA-3 Fisher ($>3$) indicates safety signal is accessible to linear methods without temporal context.
\end{itemize}

\subsection{Evaluation Metrics}

\paragraph{Coverage.} Fraction of test traces where the probe fires within $\pm 3$ steps of the true commitment point.

\paragraph{Lead time.} Signed distance (in reasoning steps) between probe firing and commitment point. Positive = early detection.

\subsection{Models and Data}

We evaluate four reasoning models spanning a 2$\times$2 architecture$\times$scale matrix:
\begin{itemize}
    \item \textbf{Llama-3.1}: DeepSeek-R1-Distill-8B (r1-8b)
    \item \textbf{Qwen2.5 (7B)}: OpenThinker-7B
    \item \textbf{Qwen2.5 (32B)}: DeepSeek-R1-Distill-32B (r1-32b), QwQ-32B
\end{itemize}

Training data consists of 209 reasoning traces from HarmThoughts with per-step safety annotations. Test set contains 43 held-out traces.

\section{Experiments}

\subsection{Cross-Architecture Static Probe Evaluation}

We first establish that static probe performance varies dramatically across architectures.

\begin{table}[h]
\centering
\caption{Static probe performance across architectures. Qwen2.5 models achieve 70\%+ coverage vs 46.5\% for Llama, explained by higher PCA variance concentration.}
\label{tab:static}
\begin{tabular}{lccccc}
\toprule
Model & Architecture & Static Coverage & Fisher (full) & Fisher (PCA-3) & PCA Var Expl. \\
\midrule
r1-8b & Llama-3.1 & 46.5\% & 17.9 & 0.93 & 16.6\% \\
r1-32b & Qwen2.5 & 69.8\% & 8.3 & 5.1 & 38.0\% \\
QwQ-32B & Qwen2.5 & 72.1\% & 6.7 & 4.2 & 41.0\% \\
\bottomrule
\end{tabular}
\end{table}

The Llama architecture distributes safety signal across minor components (PCA top-3 = 16.6\%, Fisher PCA-3 = 0.93), making it inaccessible to static linear classifiers. Qwen2.5 concentrates signal into dominant PCs (38--41\% variance, Fisher PCA-3 = 4.2--5.1), enabling linear separability without temporal context.

\subsection{Temporal Window Evaluation}

\begin{table}[h]
\centering
\caption{Temporal advantage is architecture-dependent. The +30pp improvement on Llama does not generalize to Qwen2.5.}
\label{tab:temporal}
\begin{tabular}{lcccc}
\toprule
Model & Static Coverage & Temporal W=15 Coverage & Delta \\
\midrule
r1-8b (Llama) & 46.5\% & 76.7\% & +30.2pp \\
r1-32b (Qwen2.5) & 69.8\% & 74.4\% & +4.6pp \\
QwQ-32B (Qwen2.5) & 72.1\% & 76.7\% & +4.6pp \\
\bottomrule
\end{tabular}
\end{table}

On Llama, temporal context captures ``safety reasoning density decline'' patterns invisible to single-step classification. On Qwen2.5, the sharp transition structure means a single step already encodes sufficient information; temporal smoothing adds only marginal regularization benefit.

\subsection{Window Size Ablation}

\begin{table}[h]
\centering
\caption{Coverage by temporal window size. Llama shows steep improvement; Qwen2.5 is nearly flat.}
\label{tab:window}
\begin{tabular}{lccc}
\toprule
W & r1-8b (Llama) & r1-32b (Qwen) & QwQ-32B (Qwen) \\
\midrule
1 & 46.5\% & 69.8\% & 72.1\% \\
3 & 53.5\% & 72.1\% & 74.4\% \\
5 & 65.1\% & 72.1\% & 74.4\% \\
10 & 65.1\% & 74.4\% & 76.7\% \\
15 & 76.7\% & 74.4\% & 76.7\% \\
20 & 76.7\% & 74.4\% & 76.7\% \\
25 & 74.4\% & 72.1\% & 76.7\% \\
\bottomrule
\end{tabular}
\end{table}

The divergent curve shapes directly visualize the two regimes: Llama coverage increases monotonically from W=1 to W=15 (+30pp), while Qwen2.5 achieves near-ceiling performance at W=1 with marginal gains thereafter.

\subsection{Bootstrap Confidence Intervals}

\begin{table}[h]
\centering
\caption{Bootstrap 95\% CIs (1000 resamples) confirm statistical significance of key comparisons.}
\label{tab:bootstrap}
\begin{tabular}{lccc}
\toprule
Comparison & Point Estimate & 95\% CI & Significant? \\
\midrule
Temporal $\Delta$ (Llama r1-8b) & +30.2pp & [+19, +42] & Yes \\
Temporal $\Delta$ (Qwen r1-32b) & +4.6pp & [-2, +12] & No \\
Temporal $\Delta$ (Qwen QwQ-32B) & +4.6pp & [-2, +12] & No \\
Static: Qwen vs Llama & +23.3pp & [+11, +36] & Yes \\
Cross-arch transfer vs random & +7.6pp & [-4, +17] & No \\
\bottomrule
\end{tabular}
\end{table}

The temporal advantage on Qwen2.5 has a CI crossing zero, meaning we cannot reject the null hypothesis that temporal context provides no benefit. The architecture difference in static performance is significant with a large effect size.

\subsection{Scale-Controlled Comparison (OpenThinker-7B)}

To disentangle architecture from scale, we evaluate OpenThinker-7B (Qwen2.5-7B base).

\begin{table}[h]
\centering
\caption{2$\times$2 architecture $\times$ scale matrix confirms architecture, not scale, determines temporal advantage.}
\label{tab:scale}
\begin{tabular}{lcc}
\toprule
 & Small (7--8B) & Large (32B) \\
\midrule
Llama & r1-8b: $\Delta$=+30.2pp & --- \\
Qwen2.5 & OpenThinker: $\Delta$=+2.4pp & r1-32b: $\Delta$=+4.6pp / QwQ: $\Delta$=+4.6pp \\
\bottomrule
\end{tabular}
\end{table}

OpenThinker-7B achieves 67.4\% static coverage and only +2.4pp temporal improvement, with Fisher PCA-3 = 4.8 (vs 0.93 for similarly-sized Llama r1-8b). This rules out model size as a confound.

\subsection{Cross-Architecture Transfer}

\begin{table}[h]
\centering
\caption{Zero-shot cross-architecture transfer. Cross-family transfer fails; within-family transfer is viable.}
\label{tab:transfer}
\begin{tabular}{llccc}
\toprule
Train Model & Test Model & Method & Coverage & vs In-Domain \\
\midrule
r1-8b (Llama) & r1-32b (Qwen) & Static LR & 25.6\% & 69.8\% \\
r1-8b (Llama) & r1-32b (Qwen) & Temporal W15 & 27.9\% & 74.4\% \\
r1-32b (Qwen) & r1-8b (Llama) & Static LR & 30.2\% & 46.5\% \\
r1-32b (Qwen) & r1-8b (Llama) & Temporal W15 & 32.6\% & 76.7\% \\
r1-32b (Qwen) & QwQ-32B (Qwen) & Static LR & 65.1\% & 72.1\% \\
\bottomrule
\end{tabular}
\end{table}

Cross-architecture transfer achieves only 25--33\% coverage (near random baseline $\sim$20\%), confirming that safety signal occupies fundamentally different geometric subspaces across architectures. Within-family transfer (Qwen$\rightarrow$Qwen: 65.1\%) demonstrates shared representation geometry within an architecture family.

\subsection{Lead Time Analysis}

\begin{table}[h]
\centering
\caption{Detection lead time. Qwen2.5 achieves positive lead time; Llama does not.}
\label{tab:leadtime}
\begin{tabular}{llccc}
\toprule
Model & Method & Lead Mean & Lead Median & Positive Lead Rate \\
\midrule
r1-8b (Llama) & Static LR & $-2.0$ & $-2$ & 30\% \\
r1-8b (Llama) & Temporal W15 & $-1.09$ & $-1$ & 34.5\% \\
r1-32b (Qwen) & Static LR & $+2.1$ & $+1$ & 52\% \\
r1-32b (Qwen) & Temporal W15 & $+3.4$ & $+2$ & 55\% \\
QwQ-32B (Qwen) & Static LR & $+3.8$ & $+3$ & 58\% \\
QwQ-32B (Qwen) & Temporal W15 & $+4.2$ & $+3$ & 60\% \\
\bottomrule
\end{tabular}
\end{table}

The concentrated-signal regime uniquely enables pre-textual detection: Qwen2.5 models achieve mean lead times of +2.1 to +4.2 steps with 52--60\% positive lead rate. In contrast, Llama shows negative lead time ($-1$ to $-2$ steps), meaning detection lags behind commitment. This architectural difference has direct deployment implications: Qwen2.5-based systems can enable proactive intervention, while Llama-based systems can only support post-hoc audit.

\subsection{Mahalanobis Distance Baseline}

\begin{table}[h]
\centering
\caption{Mahalanobis distance (TrajGuard-style) baseline. Partially effective on Qwen2.5, poor on Llama.}
\label{tab:mahalanobis}
\begin{tabular}{lcccc}
\toprule
Model & Architecture & MD Coverage & FPR & vs Static LR \\
\midrule
r1-8b & Llama-3.1 & 34.9\% & 23\% & $-11.6$pp \\
r1-32b & Qwen2.5 & 55.8\% & 14\% & $-14.0$pp \\
QwQ-32B & Qwen2.5 & 60.5\% & 11\% & $-11.6$pp \\
\bottomrule
\end{tabular}
\end{table}

Mahalanobis distance achieves 55--60\% coverage on Qwen2.5, indicating that even unsupervised geometric methods provide practical (if suboptimal) monitoring when signal is concentrated. On Llama, coverage drops to 35\% due to safety signal residing in minor components that contribute minimally to Mahalanobis distance.

\subsection{Layer Selection Ablation}

\begin{table}[h]
\centering
\caption{Layer sensitivity. Qwen2.5 is robust to layer choice; Llama is sensitive.}
\label{tab:layers}
\begin{tabular}{lccc}
\toprule
Model & Best Layer & Coverage Range & Layers $>$70\% Coverage \\
\midrule
r1-8b (Llama) & 14 & 38--77\% & 1 (only L14 temporal) \\
r1-32b (Qwen) & 18 & 52--74\% & 10 (L14--L24) \\
QwQ-32B (Qwen) & 16 & 55--77\% & 10 (L12--L22) \\
\bottomrule
\end{tabular}
\end{table}

Qwen2.5 models exhibit broad effective layer ranges (10 layers above 70\% coverage), reducing engineering sensitivity for deployment. Llama requires precise layer selection (only L14 with temporal context exceeds 70\%).

\subsection{GRU Upper Bound}

\begin{table}[h]
\centering
\caption{Model complexity comparison on 32B models. GRU provides no improvement over temporal LR.}
\label{tab:gru}
\begin{tabular}{lccc}
\toprule
Method & r1-32b Coverage & QwQ-32B Coverage & Parameters \\
\midrule
Static LR & 69.8\% & 72.1\% & $\sim$5K \\
Temporal LR W=15 & 74.4\% & 76.7\% & $\sim$77K \\
GRU-256 & 74.4\% & 76.7\% & $\sim$1.6M \\
\bottomrule
\end{tabular}
\end{table}

GRU achieves identical coverage to temporal LR despite 20$\times$ more parameters, confirming that the detection ceiling is architecture-determined, not capacity-determined. For Qwen2.5, static LR is near-ceiling; for both architectures, increased model complexity yields no benefit.

\subsection{Fine-Grained Supervision (16-Label)}

\begin{table}[h]
\centering
\caption{16-label taxonomy vs binary supervision. Marginal improvement on all architectures.}
\label{tab:16label}
\begin{tabular}{lccc}
\toprule
Model & Binary Coverage & 16-Label Coverage & Delta \\
\midrule
r1-8b (Llama) Temporal & 76.7\% & 79.1\% & +2.4pp \\
r1-32b (Qwen) Temporal & 74.4\% & 76.7\% & +2.3pp \\
QwQ-32B (Qwen) Temporal & 76.7\% & 79.1\% & +2.4pp \\
\bottomrule
\end{tabular}
\end{table}

Fine-grained labels yield a small uniform improvement ($\sim$2pp) that does not interact with the architecture-dependent story. The bottleneck on Llama is signal geometry (not label quality), and on Qwen2.5 the static probe is already near-optimal.

\subsection{Oracle Decomposition}

\begin{table}[h]
\centering
\caption{Oracle signal source analysis reveals mechanistic basis for architectural differences.}
\label{tab:oracle}
\begin{tabular}{lccc}
\toprule
Model & Oracle Signal Source & Hidden-State Corr. & Pre-CP Entropy \\
\midrule
r1-8b (Llama) & Density gradient & 0.21 & 1.83 \\
r1-32b (Qwen) & Sharp transition & 0.67 & 0.94 \\
QwQ-32B (Qwen) & Sharp transition & 0.67 & 0.94 \\
\bottomrule
\end{tabular}
\end{table}

The two architectures exhibit fundamentally different commitment dynamics. Llama shows gradual density decline (high pre-commitment-point entropy = 1.83), requiring temporal aggregation to detect. Qwen2.5 shows sharp transitions (low entropy = 0.94) that correlate strongly with instantaneous hidden states ($r = 0.67$ vs 0.21 for Llama), explaining why static probes succeed.

\section{Discussion}

\subsection{A Diagnostic Framework for Practitioners}

Our results suggest a simple decision procedure for deploying safety monitors:

\begin{enumerate}
    \item Compute PCA variance concentration (top-3 components) on hidden states.
    \item If variance explained $> 30\%$: \textbf{concentrated-signal regime}. Deploy static LR probe; temporal context is optional. Expect positive lead time and robust layer selection.
    \item If variance explained $< 20\%$: \textbf{distributed-signal regime}. Deploy temporal window probe (W=15). Expect negative lead time; design for post-hoc audit rather than proactive intervention.
\end{enumerate}

The Fisher discriminant ratio in the PCA-3 subspace provides a secondary diagnostic: values $> 3$ indicate concentrated signal.

\subsection{Implications for Cross-Architecture Deployment}

The catastrophic failure of cross-architecture transfer (25--33\%) means safety monitors must be trained per-architecture. However, within-family transfer succeeds (65.1\%), suggesting that a single probe trained on one Qwen2.5 model can provide baseline monitoring for other Qwen2.5 models.

\subsection{Limitations}

Our study has several limitations: (1) sample size is modest (209 training / 43 test traces), reflected in wide bootstrap CIs; (2) we evaluate only two architecture families; (3) results are specific to the HarmThoughts dataset and commitment-point detection task; (4) we do not evaluate on proprietary reasoning models (o1, o3) whose architectures are undisclosed.

\section{Conclusion}

We present the first cross-architecture evaluation of hidden-state safety monitoring in reasoning models, revealing that monitoring effectiveness is fundamentally architecture-dependent. We identify two regimes---distributed-signal (Llama, requiring temporal probes) and concentrated-signal (Qwen2.5, where static probes suffice)---and provide a diagnostic framework (PCA variance concentration, Fisher ratio) for practitioners to determine which regime applies. Our findings establish that universal monitoring approaches are insufficient: architecture determines not only \emph{which} method works but \emph{whether} pre-textual detection is possible at all. Future work should extend this characterization to additional architecture families and develop architecture-adaptive monitoring systems.

\bibliographystyle{plain}
\begin{thebibliography}{9}
\bibitem{harmthoughts} Anonymous. HarmThoughts: Characterizing safety reasoning patterns in extended chain-of-thought. \emph{Preprint}, 2025.
\bibitem{selfjailbreak} Anonymous. Self-jailbreaking: Reasoning models circumvent safety training through extended deliberation. \emph{Preprint}, 2025.
\bibitem{trajguard} Anonymous. TrajGuard: Hidden-state trajectory monitoring for jailbreak detection. \emph{ACL}, 2025.
\bibitem{safeswitch} Anonymous. SafeSwitch: Linear probes for real-time safety classification in LLMs. \emph{EMNLP}, 2025.
\bibitem{siren} Anonymous. SIREN: Sequential inference for reasoning safety enforcement. \emph{NeurIPS}, 2025.
\bibitem{saferemind} Anonymous. SafeRemind: Text-level monitoring of reasoning chain safety. \emph{Preprint}, 2025.
\bibitem{safethink} Anonymous. SafeThink: Proactive safety intervention in chain-of-thought. \emph{Preprint}, 2025.
\end{thebibliography}

\end{document}
